\subsection{NORMS AND TRACES RELATIVE TO A MODULE}

DEFINITION 1. Let M be an A-module admitting a finite basis as a K-module. For all \(a \in \mathbf{A}\) , let \(a_{\mathbf{M}}\) be the endomorphism \(x \mapsto ax\) of the K-module M. The trace, determinant and characteristic polynomial of \(a_{\mathbf{M}}\) are called respectively the trace, norm and characteristic polynomial of \(z\) relative to M.

The trace and norm of a are therefore elements of K, denoted respectively by \(\mathrm{Tr}_{\mathbf{M}/\mathbf{K}}(a)\) and \(\mathrm{N}_{\mathbf{M}/\mathbf{K}}(a)\) ; the characteristic polynomial of a is an element of K{[}X{]}, denoted by \(\mathrm{Pc}_{\mathbf{M}/\mathbf{K}}(a;\mathbf{X})\) . We omit K in the above notation when there is no risk of confusion.

From the properties of the trace and determinant of an endomorphism (II, § 4, no. 3 and § 8, no. 1) we obtain the relations

\[
\operatorname{Tr} _ {\mathbf {M}} (a + a ^ {\prime}) = \operatorname{Tr} _ {\mathbf {M}} (a) + \operatorname{Tr} _ {\mathbf {M}} (a ^ {\prime})\tag{1}
\]

\begin{enumerate}
\def\labelenumi{(\arabic{enumi})}
\setcounter{enumi}{1}
\tightlist
\item
\end{enumerate}

\[
\operatorname{Tr} _ {\mathbf {M}} (a a ^ {\prime}) = \operatorname{Tr} _ {\mathbf {M}} (a ^ {\prime} a)\tag{3}
\]

\[
\mathrm{N} _ {\mathbf {M}} (a a ^ {\prime}) = \mathrm{N} _ {\mathbf {M}} (a) \mathrm{N} _ {\mathbf {M}} (a ^ {\prime})
\]

for all \(a, a'\) in A.

Let \((e_i)_{1\leq i\leq n}\) be a basis of the K-module M and \((m_{ij}(a))\) the matrix of the endomorphism \(a_{\mathbf{M}}\) with respect to this basis. The functions \(m_{ij}\) are linear forms on the K-module A and

\begin{enumerate}
\def\labelenumi{(\arabic{enumi})}
\setcounter{enumi}{3}
\tightlist
\item
\end{enumerate}

\[
\operatorname{Tr} _ {\mathbf {M}} (a) = \sum_ {i = 1} ^ {n} m _ {i i} (a)\tag{5}
\]

\[
\mathrm{N} _ {\mathbf {M}} (a) = \det (m _ {i j} (a))\tag{6}
\]

\[
\operatorname{Pc} (a; \mathbf {X}) = \det (\delta_ {i j} \mathbf {X} - m _ {i j} (a)).
\]

It follows from the method of calculating a determinant (§ 8, no. 11, formula (50)) that

\[
\mathrm{Pc} _ {\mathbf {M}} (a; \mathbf {X}) = \mathbf {X} ^ {n} + c _ {1} \mathbf {X} ^ {n - 1} + \dots + c _ {n}\tag{7}
\]

where

\[
c _ {1} = - \operatorname{Tr} _ {\mathbf {M}} (a), \quad c _ {n} = (- 1) ^ {n} \mathrm{N} _ {\mathbf {M}} (a).\tag{8}
\]

For \(\lambda \in \mathbf{K}\) ,

\[
\operatorname{Tr} _ {\mathbf {M}} (\lambda) = n. \lambda , \quad \mathrm{N} _ {\mathbf {M}} (\lambda) = \lambda^ {n}, \quad \operatorname{Pc} _ {\mathbf {M}} (\lambda ; \mathbf {X}) = (\mathbf {X} - \lambda) ^ {n}.\tag{9}
\]

Let \(\mathbf{K}'\) be a commutative K-algebra. Let \(\mathbf{M}' = \mathbf{K}' \otimes_{\mathbf{K}} \mathbf{M}\) and \(\mathbf{A}' = \mathbf{K}' \otimes_{\mathbf{K}} \mathbf{A}\) , so that \(\mathbf{M}'\) has an A'-module structure (§ 4, Example 2). As a \(\mathbf{K}'\) -module \(\mathbf{M}'\) admits the basis consisting of the \(1 \otimes e_i\) ( \(1 \leqslant i \leqslant n\) ) and the matrix of \(a_{\mathbf{M}}\) with respect to \((e_i)\) is equal to the matrix of \((1 \otimes a)_{\mathbf{M}'}\) with respect to \((e_i')\) . Then

\[
\operatorname{Tr} _ {\mathbf {M} ^ {\prime}} (1 \otimes a) = \operatorname{Tr} _ {\mathbf {M}} (a). 1, \quad \mathrm{N} _ {\mathbf {M} ^ {\prime}} (1 \otimes a) = \mathrm{N} _ {\mathbf {M}} (a). 1\tag{12}
\]

\[
\mathrm{Pc} _ {\mathbf {M} ^ {\prime}} (1 \otimes a; \mathbf {X}) = \mathrm{Pc} _ {\mathbf {M}} (a; \mathbf {X}). 1
\]

for all \(a \in \mathbf{A}\) , where 1 denotes the unit element of \(\mathbf{K}'\) . If in particular we take \(\mathbf{K}' = \mathbf{K}[\mathbf{X}]\) , then

\[
\mathrm{Pc} _ {\mathbf {M} / \mathbf {K}} (a; \mathbf {X}) = \mathrm{N} _ {\mathbf {M} [ \mathbf {X} ] / \mathbf {K} [ \mathbf {X} ]} (\mathbf {X} - a).\tag{13}
\]

